\section{PR-III Finite-tier Uniqueness and Closure Classification}
\label{sec:pr3-branch-countermodels}

PR-III reports a reproducible no-fit radiative and cross-sector closure at its declared finite tier \cite{PRIII}.  That order is a defining condition of the PR class rather than a parameter varied inside it.  PR-IV classifies equation presentations at that fixed order.

\subsection{Electromagnetic Radiative Monomials}

The PR-III electromagnetic corrections use dimensionless factors drawn from $p_1$, $\Delta_{\rm EW}$, $s_W^2$, and $N_{\rm gen}^{-1/2}$, including
\begin{align}
 D_1&=-p_1\Delta_{\rm EW}s_W^2,\\
 D_2&=D_1\Delta_{\rm EW}N_{\rm gen}^{-1/2}.
\end{align}
If one drops the declared PR-III finite-tier composition grammar, one can write formal expressions such as
\begin{equation}
 -p_1\Delta_{\rm EW}s_W^2N_{\rm gen}^{-1/2},
 \qquad
 -p_1\Delta_{\rm EW}^2s_W^2,
 \qquad
 -p_1\Delta_{\rm EW}(s_W^2)^2
 \label{eq:pr35-pr3-monomial-counterexamples}
\end{equation}
from the same symbols.  The bounded-degree audit enumerated 23 such expressions through total degree six.  Twenty-one violate the inherited finite-tier composition order and therefore lie outside
$\mathfrak A_{\rm PR}^{\rm can}$; they are not alternative PR equations. Inside the declared grammar, $D_1$ and $D_2$ are the sole two normal forms and carry no adjustable exponent or coefficient.

\subsection{Charged and Neutral Electroweak Coefficients}

The $C_2$ correction uses the charged/neutral coefficient pair $(2,-1)$ and the neutral $D_3$ correction uses coefficient $2$.  Before any representation-trace theorem is imposed, the preregistered sign gate alone is satisfied, for example, by $(1,-1)$, $(2,-1)$, and $(3,-2)$.  Likewise the positive neutral coefficients $1,2,3,4$ all satisfy a strict positivity-only gate.
The small-integer audit found 16 opposite-sign $C_2$ pairs and four positive $D_3$ coefficients in its stated range.  More strongly, the exact real-symmetric gauge-index coefficient classification at the fixed $D_3$ Lorentz/momentum kernel in Appendix~\ref{app:d3-coefficient-rigidity} shows that the three linear structural gates leave $D_3(k)=\rho kP_Z$ for every real $k$ within that matrix class.  Adding the separate strict positivity-only gate restricts this to $k>0$, still an uncountable half-line.  Fixed rational members are discrete, contain no fitted continuous coefficient, and are not
convention-equivalent once kinetic and field normalization are frozen.

The PR-III pair and coefficient are therefore fixed by a framework-native projector-trace and response closure, not by the sign gate or diagnostic residual.  Appendix \ref{app:d3-coefficient-rigidity} now proves that the canonical charged-adjoint contraction has neutral representation trace two and that the four canonical defects vanish only for $D_3=2\rho P_Z$. Theorem~\ref{thm:pr35-unique-quadratic-response} proves that the normalized charged-adjoint functional is the unique minimal response and therefore fixes $k=2$ within the declared PR class.  A determinant-Hessian expression satisfying the same normalized physical-$Z$ map is an equivalent microscopic representation of this result.  The all-real family obtained after omitting canonical normalization lies outside the PR class.  The locked $M_Z$ gate permits an interval containing $k=1,2,3$, confirming that the exact normalized-response defect, rather than residual ranking, supplies the selection.  The $C_1$ transport rule, $C_2$ trace pair, and multiplicative neutrino envelope are fixed inherited finite-tier maps; changing their operator grammar defines a different theory rather than another solution of the declared PR equations.

The strong-sector one-loop coefficient is fixed once its representation and active-threshold inventory are imposed.  Its theorem domain is exactly the declared one-loop strong class.  Higher-order or nonperturbative operators lie outside that finite-tier class and are not additional solutions of it.
